\subsection{Profundidad de Anidamiento}

\subsubsection{Regla: máximo tres niveles de anidamiento}

\textbf{Regla:} No se permitirá anidar más de tres niveles de profundidad de estructuras de control (\texttt{if}, \texttt{for}, \texttt{try}) dentro de un mismo método. Cuando el algoritmo lo requiera, las condiciones lógicas deberán combinarse en una sola línea o el fragmento deberá extraerse a un método privado.

\textbf{Descripción:} McConnell señala que la profundidad de anidamiento excesiva es una de las causas principales de la complejidad estructural de una rutina y recomienda mantenerla baja para conservar la legibilidad del código (McConnell, 2004).

\textbf{Código correcto:}
\begin{lstlisting}[style=csharpstyle]
public void ProcessTurn(Player player)
{
    if (player is not null && player.IsActive)
    {
        ExecuteTurn(player);
    }
}

private void ExecuteTurn(Player player)
{
    try
    {
        _board.MovePlayer(player);
    }
    catch (InvalidMoveException ex)
    {
        _logger.Warning($"Invalid move for {player.Name}: {ex.Message}");
    }
}
\end{lstlisting}

\textbf{Código incorrecto:}
\begin{lstlisting}[style=csharpstyle]
public void ProcessTurn(Player player)
{
    if (player is not null)
    {
        if (player.IsActive)
        {
            try
            {
                if (_board.HasSpaceFor(player))
                {
                    _board.MovePlayer(player);
                }
            }
            catch (InvalidMoveException ex)
            {
                _logger.Warning($"Invalid move for {player.Name}: {ex.Message}");
            }
        }
    }
}
\end{lstlisting}
